% Corrigé intégral imprimable du dernier TD, avec planches UML explicatives.
% Les solutions sont proposées et argumentées; elles ne sont pas le corrigé officiel.
\documentclass[10pt,a4paper]{article}
\usepackage{fontspec}
\setmainfont{Carlito}
\setsansfont{Carlito}
\usepackage[french]{babel}
\usepackage[a4paper,margin=17mm,headheight=17pt,footskip=12mm]{geometry}
\usepackage{microtype,xcolor,hyperref,booktabs,tabularx,longtable,array,multicol,enumitem}
\usepackage{tikz}
\usetikzlibrary{positioning,arrows.meta,fit,calc,shapes.geometric,backgrounds}
\usepackage[most]{tcolorbox}
\usepackage{titlesec,fancyhdr,lastpage}
\definecolor{navy}{HTML}{153F56}\definecolor{teal}{HTML}{157A86}
\definecolor{pale}{HTML}{EDF6F7}\definecolor{warm}{HTML}{FFF4E4}\definecolor{bodygray}{HTML}{445866}
\hypersetup{colorlinks=true,linkcolor=teal,urlcolor=teal,pdftitle={Corrige integral - Organisation de courses de fond}}
\renewcommand{\arraystretch}{1.20}
\setlist[itemize]{leftmargin=*,itemsep=1.5pt,topsep=3pt}
\setlist[enumerate]{leftmargin=*,itemsep=1.5pt,topsep=3pt}
\tcbset{before skip=7pt,after skip=7pt,colback=pale,colframe=teal!40,arc=2mm,boxrule=.6pt,left=8pt,right=8pt,top=6pt,bottom=6pt}
\newtcolorbox{memo}[1][]{colback=pale,colframe=teal,title={À retenir},fonttitle=\bfseries,#1}
\newtcolorbox{warningbox}[1][]{colback=warm,colframe=orange!55!black,title={Piège à éviter},fonttitle=\bfseries,#1}
\newtcolorbox{exobox}[1][]{colback=white,colframe=navy!40,title={Exemple difficile},fonttitle=\bfseries,#1}
\titleformat{\section}{\Large\bfseries\color{navy}}{\thesection.}{.6em}{}
\titleformat{\subsection}{\large\bfseries\color{teal}}{\thesubsection}{.6em}{}
\titleformat{\subsubsection}{\normalsize\bfseries\color{navy}}{\thesubsubsection}{.6em}{}
\titlespacing*{\section}{0pt}{17pt}{8pt}
\titlespacing*{\subsection}{0pt}{12pt}{5pt}
\pagestyle{fancy}\fancyhf{}
\fancyhead[L]{\footnotesize\color{bodygray} M2 SIA -- Organisation de courses de fond}
\fancyhead[R]{\footnotesize\color{bodygray} Corrigé commenté}
\fancyfoot[L]{\footnotesize\color{bodygray} Support d'entraînement, non officiel -- octobre 2026}
\fancyfoot[R]{\footnotesize\color{bodygray} Page \thepage/\pageref{LastPage}}
\renewcommand{\headrulewidth}{.35pt}
\color{navy}\setlength{\parindent}{0pt}\setlength{\parskip}{5pt}
\newcommand{\source}[1]{\par{\footnotesize\color{bodygray}\textit{Source : #1}}}
\tikzset{actor/.style={draw=navy,rounded corners=1pt,minimum width=16mm,minimum height=8mm,fill=white,font=\scriptsize,align=center},usecase/.style={draw=teal,ellipse,minimum width=27mm,minimum height=11mm,fill=pale,font=\scriptsize,align=center},cbox/.style={draw=teal,rounded corners=2pt,fill=pale,minimum width=26mm,minimum height=9mm,align=center,font=\scriptsize},dbox/.style={draw=teal,rectangle,fill=pale,minimum width=24mm,minimum height=8mm,align=center,font=\scriptsize},assoc/.style={-Stealth,draw=navy,thick},umlarrow/.style={-Stealth,draw=navy,dashed},life/.style={draw=bodygray,densely dashed},state/.style={draw=teal,rounded corners=4pt,fill=pale,minimum width=28mm,minimum height=11mm,align=center,font=\scriptsize}}
\begin{document}
\thispagestyle{empty}
\vspace*{22mm}
{\color{teal}\rule{35mm}{3pt}}\\[12mm]
{\Huge\bfseries Corrigé intégral\par}
\vspace{4mm}{\LARGE Organisation de courses de fond\par}
\vspace{4mm}{\large Les 9 questions du TD final + diagrammes détaillés\par}
\vspace{14mm}
\begin{tcolorbox}[colback=pale]
\textbf{But :} disposer d'un modèle de réponse complet, simple à lire et assez difficile pour s'entraîner à l'examen.\\
\textbf{Contenu :} cas d'utilisation, description textuelle, séquences N1, classes métier et paquetages, états, deux DCP, communication, séquence N3.\\
\textbf{Méthode :} textes du TD (6 pages) + fiches méthodologiques ASI C1--C5. Hypothèses indiquées quand le TD ne décide pas lui-même.\end{tcolorbox}
\vfill
{\small\color{bodygray} Document d'entraînement indépendant. Ne pas confondre avec la copie envoyée à M. Lahmer le 21 octobre 2026, ni avec un corrigé validé par l'enseignant.}
\newpage
\tableofcontents
\newpage
\section*{Mode d'emploi avant d'ouvrir les diagrammes}
\addcontentsline{toc}{section}{Mode d'emploi}
\begin{memo}
À l'examen, vous avez un long énoncé : soulignez d'abord les acteurs, les événements, les données métier, les délais et les exceptions. Le TD contient trois canaux (internet, guichet, courrier) et une variante de club (triathlon et vélo). \textbf{Ne fusionnez pas les règles de ces quatre scénarios.}
\end{memo}
\begin{tabularx}{\linewidth}{@{}p{40mm}X@{}}\toprule
\textbf{Mot du sujet} & \textbf{Effet sur la modélisation} \\ \midrule
\textit{« au plus tard »} & Contrainte de date sur le cas et garde sur la transition ; pas nécessairement minuterie UML.\\
\textit{« si... sinon... »} & Branche alternative \texttt{alt} dans une séquence, garde \texttt{[condition]} dans les états.\\
\textit{« a déjà participé »} & Réutiliser les données de \texttt{Concurrent}, mais vérifier le certificat à nouveau.\\
\textit{« club uniquement »} & Acteur Club distinct ; formulaire traité pour chaque membre ; licence à la place du certificat.\\
\textit{« cinq meilleurs temps »} & Calcul métier du classement de club ; au plus quatre finisseurs signifie disqualification.\\
\textit{« transpondeur »} & Entité physique liée au dossard, plusieurs passages et restitution obligatoire.\\
\bottomrule\end{tabularx}
\textbf{Hypothèses de modélisation :} H1 : un agent enregistre au guichet/courrier ; H2 : la banque et le service SMS sont externes ; H3 : \texttt{Inscription} possède un statut commun ; H4 : une réservation de dossard et une affectation de puce sont gérées par des objets métier, sans inventer un calendrier précis de distribution ; H5 : le système peut enregistrer des résultats de vérification. Tout choix non explicite est annoncé comme hypothèse.
\newpage
% Corrigé pédagogique original, fondé sur « Organisation de courses de fond » (6 p.).
% Aucun diagramme du professeur reproduit ; choix de modélisation justifiés.
\section{Cas final : organisation de courses de fond -- correction complète des 9 travaux}
\begin{memo}
Cette partie est la \textbf{référence de travail prioritaire}. Elle répond \textbf{aux neuf demandes} du dernier TD. Le texte du professeur fait autorité. Ce corrigé est \textbf{une proposition argumentée}, pas une correction officielle. Il faut déclarer nos hypothèses H1--H4 quand les règles ne sont pas explicites.
\end{memo}
\textbf{Faits à retenir :} course libre dès 16 ans, semi-marathon dès 18 ans, marathon dès 20 ans ; triathlon et vélo : inscription par club, en ligne ou par courrier, avec licence au lieu de certificat. Préinscription internet jusqu'à J--15, sous réserve de places. Sur place : dossier terminé au moins 1 h avant le départ, pas de triathlon/vélo. Par courrier : terminé au moins J--5. Paiement invalide ou certificat non conforme : annulation en ligne ; par courrier, demande de document et maintien à l'étape adaptée. Club classé par ses cinq meilleurs temps, sinon disqualifié (au plus quatre arrivants). SMS en option payante à 1 euro. Transpondeur : numéro de dossard, temps et passages, restitution obligatoire.

\begin{warningbox}
\textbf{Hypothèses déclarées :} H1 : un agent du secrétariat saisit les inscriptions reçues sur papier ou au guichet ; H2 : on sépare l'état de l'\texttt{Inscription} du support de paiement ; H3 : le statut d'inscription suit des étapes communes, sauf les parcours plus rapides qui sautent plusieurs étapes ; H4 : la catégorie d'âge et l'admissibilité sont vérifiées par \texttt{Epreuve} (le texte n'impose aucune classe responsable). Ne transformez pas ces choix en « faits du sujet ».
\end{warningbox}

\subsection{1. Cas d'utilisation : acteurs, objectifs et paquetages}
\textbf{Acteurs primaires :} Concurrent, Club (pour vélo et triathlon), AgentSecrétariat, ServiceCourses, Administrateur (si gestion technique précisée). \textbf{Acteurs secondaires :} Banque/PrestataireCarte, SystèmeChronométrage (antennes et puces externes, sous H5 si le périmètre du logiciel ne les inclut pas), ServiceSMS (si externe). Un transpondeur passif n'est \textbf{pas automatiquement} acteur ; seul le dispositif qui échange effectivement des messages avec le système peut l'être.
\begin{center}
\begin{tikzpicture}[scale=.84,transform shape,node distance=9mm and 17mm]
\node[actor] (cc) at (0,2.8) {Concurrent};\node[actor] (club) at (0,.8) {Club};
\node[actor] (sec) at (0,-1.3) {Secrétariat};\node[actor] (srv) at (0,-3.3) {Service courses};
\node[usecase] (online) at (4.5,3.2) {S'inscrire sur internet};
\node[usecase] (cour) at (4.5,1.7) {Inscrire par courrier};
\node[usecase] (place) at (4.5,.2) {Inscrire sur place};
\node[usecase] (consult) at (4.5,-1.3) {Rechercher un classement};
\node[usecase] (classer) at (4.5,-2.8) {Enregistrer classement};
\node[usecase] (result) at (4.5,-4.3) {Enregistrer les temps};
\node[actor] (bank) at (8.7,3.4) {Banque};
\node[actor] (chrono) at (8.7,-4.3) {Système de\\chronométrage};
\draw[navy] (2.7,-5.05) rectangle (7.1,4.05);
\node[font=\scriptsize,fill=white] at (4.9,4.04){Système d'organisation des courses};
\draw[navy] (cc)--(online);\draw[navy](cc)--(consult);
\draw[navy](club)--(online);\draw[navy](club)--(cour);
\draw[navy](sec)--(cour);\draw[navy](sec)--(place);
\draw[navy](srv)--(classer);\draw[navy](srv)--(result);
\draw[navy](bank)--(online);\draw[navy](chrono)--(result);
\end{tikzpicture}
\end{center}
\textbf{Paquetages fonctionnels :} \texttt{Inscriptions}, \texttt{EpreuvesEtClubs}, \texttt{Classements}, \texttt{Chronometrage}. Un paquetage de cas n'est pas une étape du scénario. Les cas « consulter conditions » et « télécharger formulaire » appartiennent à l'espace public ; « gérer récompenses » au ServiceCourses. \textbf{Attention :} un Club n'est pas un « Concurrent amélioré » ; c'est un rôle institutionnel différent.

\subsection{2. Description textuelle détaillée : inscription sur internet}
\begin{memo}
\textbf{Méthode exacte d'ASI\_C1, page 36 :} Identification (nom, but, acteurs) ; Séquencement (déclencheur, préconditions, enchaînement nominal, postconditions) ; variantes et exceptions ; éventuellement exigences non fonctionnelles et IHM. Le TD de course de fond demande explicitement ce travail : ne pas le réduire à un diagramme.
\end{memo}
\begin{tabularx}{\linewidth}{@{}p{40mm}X@{}}
\toprule
\textbf{Nom du cas} & S'inscrire à une course sur internet.\\
\textbf{But} & Obtenir une inscription définitive après paiement accepté et vérification des pièces.\\
\textbf{Acteur principal} & Concurrent pour une course individuelle ; Club de la discipline pour triathlon ou vélo.\\
\textbf{Acteurs secondaires} & Banque (carte) ; Secrétariat (chèque, certificat ou licence).\\
\textbf{Déclencheur} & Le concurrent choisit une course et demande son inscription depuis le site web.\\
\textbf{Préconditions} & Une épreuve est publiée et accessible via le site. On \emph{contrôle pendant le scénario} la date, les places et l'admissibilité.\\
\textbf{Postcondition succès} & Une \texttt{Inscription} est \texttt{DEFINITIVE} ; paiement accepté et justificatif conforme.\\
\textbf{Postconditions sans succès} & Avant la préinscription, aucun dossier si quota/date/âge refusé. Après paiement invalide ou pièce non conforme : \texttt{ANNULEE}, concurrent prévenu. Dossier temporairement \texttt{PREINSCRIPTION} ou \texttt{FINANCIERE} en attendant les pièces.\\
\bottomrule
\end{tabularx}
\textbf{Enchaînement nominal (course individuelle ordinaire, carte) :}
\begin{enumerate}[label=\textbf{\arabic*.}]
\item Le concurrent sélectionne la course. Le système affiche ses conditions, sa date et son tarif.
\item Le système vérifie le quota, la limite J--15 et la condition d'âge liée à la discipline.
\item Le concurrent fournit ou confirme ses données personnelles. Le système retrouve les données déjà connues, si elles existent.
\item Le système crée l'inscription au statut \texttt{PREINSCRIPTION}.
\item Le concurrent choisit de payer par carte. Le système transmet la demande de paiement à la banque.
\item La banque accepte. Le système fait passer l'inscription au statut \texttt{FINANCIERE}.
\item Le concurrent envoie par courrier un certificat médical de moins d'un an. Le secrétariat le reçoit et le contrôle.
\item Si le certificat est conforme, le système passe à \texttt{DEFINITIVE} et le concurrent est informé.
\end{enumerate}
\textbf{Enchaînements alternatifs :}
\begin{itemize}
\item \textbf{A1 [étape 3 : concurrent déjà connu] :} le système propose les anciennes informations ; le concurrent les valide ou les corrige ; retour à l'étape 4.
\item \textbf{A2 [étapes 5--6 : chèque] :} le concurrent envoie son paiement par courrier ; le secrétariat le contrôle ; le statut devient \texttt{FINANCIERE} quand le chèque est accepté. Le certificat est vérifié ensuite.
\item \textbf{A3 [triathlon / vélo] :} le \texttt{Club} inscrit son membre par internet ; on contrôle l'adhésion au club spécialisé ; la \texttt{Licence} conforme remplace le certificat médical ; il n'y a \textbf{pas} d'inscription sur place.
\end{itemize}
\textbf{Exceptions :}
\begin{itemize}
\item \textbf{E1 [étape 2 : places insuffisantes ou date J--15 dépassée] :} refus avant préinscription. La condition d'âge empêche également l'inscription aux courses individuelles non autorisées.
\item \textbf{E2 [étape 6 : paiement non valide] :} inscription annulée et concurrent informé.
\item \textbf{E3 [étape 8 : certificat non conforme] :} inscription annulée et concurrent informé.
\end{itemize}
\textbf{Règles métier et remarques :} libre \(\geq16\) ans, semi \(\geq18\), marathon \(\geq20\) ; quota et date J--15 ; certificat médical datant de moins d'un an (ou licence du club spécialisé). \textbf{N'inventez pas} un upload internet de certificat ni une double autorisation/débit : ici la validation définitive est liée à la vérification du justificatif.\par
\textbf{Contrainte non fonctionnelle possible (suggestion, non donnée mot pour mot) :} limiter l'accès aux données personnelles et médicales aux personnes autorisées.
\begin{warningbox}
Un \textbf{paiement accepté} ne signifie pas \textbf{inscription définitive}. La préinscription, le paiement et le contrôle du justificatif sont \emph{trois étapes différentes}. Un courrier incomplet, lui, n'est pas automatiquement annulé : le secrétariat relance le concurrent et conserve l'étape en cours.
\end{warningbox}
\begin{tabularx}{\linewidth}{@{}p{29mm}XXX@{}}
\toprule
Pour comparer & Internet & Sur place & Courrier \\
\midrule
Date clé & Préinscription au plus tard J--15 & Dossier complet 1 h avant départ & Dossier complet au plus tard J--5\\
Qui intervient ? & Concurrent (ou club), banque, secrétariat & Concurrent, secrétariat & Concurrent/club, secrétariat\\
Mode paiement & Carte ou chèque & Espèces, carte ou chèque & Règlement joint, souvent chèque\\
Justificatif & Certificat envoyé par courrier ; licence si club & Certificat remis et contrôlé & Certificat joint ; licence si club\\
Si pièce manque & Dossier attend tant que possible & Validation immédiate impossible & Conserver étape et demander pièce\\
Triathlon/vélo & Club uniquement & \textbf{Non admis} & Club uniquement\\
\bottomrule
\end{tabularx}
\source{Organisation de courses de fond, pages 3--4 ; ASI\_C1, page 36.}

\subsection{3. Les trois séquences système N1 d'inscription}
\textbf{Règle :} un seul participant \texttt{:SystèmeCourses}, plus acteurs externes. Aucun contrôleur, écran ou DAO. Chaque case \texttt{alt} correspond à un véritable embranchement.
\begin{center}
\begin{tikzpicture}[x=.94cm,y=.49cm,scale=.91,transform shape]
\foreach \x/\t in {0/Concurrent,3.5/SystèmeCourses,7.3/Banque,11/Secrétariat}{\node[dbox] at (\x,0){\t};\draw[life](\x,-.5)--(\x,-20.8);}
\draw[assoc](0,-1)--node[above,font=\scriptsize]{1: choisirEpreuve()}(3.5,-1);
\draw[<-,dashed](0,-2)--node[above,font=\scriptsize]{2: conditions, quota, tarif}(3.5,-2);
\draw[teal](-.6,-2.6)rectangle(4.3,-4.6);
\node[font=\tiny,anchor=west] at (-.4,-2.87){break [quota/date/âge non valides]};
\draw[<-,dashed](0,-3.73)--node[above,font=\tiny]{3: refus motivé, pas de dossier}(3.5,-3.73);
\draw[assoc](0,-5.2)--node[above,font=\scriptsize]{4: saisirOuConfirmerDonnees()}(3.5,-5.2);
\draw[<-,dashed](0,-6.2)--node[above,font=\scriptsize]{5: préinscription enregistrée}(3.5,-6.2);
\draw[assoc](0,-7.2)--node[above,font=\scriptsize]{6: choisirPaiement()}(3.5,-7.2);
\draw[teal](-.6,-7.7)rectangle(11.9,-12.25);
\node[font=\tiny,anchor=west] at (-.4,-7.94){alt [carte]};
\draw[assoc](3.5,-8.75)--node[above,font=\scriptsize]{7a: demanderPaiement()}(7.3,-8.75);
\draw[<-,dashed](3.5,-9.7)--node[above,font=\scriptsize]{8a: accepté/refusé}(7.3,-9.7);
\draw[teal,densely dashed](-.6,-10.18)--(11.9,-10.18);
\node[font=\tiny,anchor=west] at (-.4,-10.4){[chèque]};
\draw[assoc](11,-11.48)--node[above,font=\scriptsize]{7b: saisirRésultatContrôleChèque()}(3.5,-11.48);
\draw[teal](-.6,-12.75)rectangle(11.9,-16.2);
\node[font=\tiny,anchor=west] at (-.4,-13){alt [paiement invalide]};
\draw[<-,dashed](0,-13.85)--node[above,font=\tiny]{9a: ANNULEE, avertissement}(3.5,-13.85);
\draw[teal,densely dashed](-.6,-14.35)--(11.9,-14.35);
\node[font=\tiny,anchor=west] at (-.4,-14.58){[paiement accepté]};
\draw[<-,dashed](0,-15.4)--node[above,font=\tiny]{9b: FINANCIERE, attendre pièce}(3.5,-15.4);
\draw[assoc](11,-16.75)--node[above,font=\scriptsize]{10: contrôlerCertificatOuLicence()}(3.5,-16.75);
\draw[teal](-.6,-17.2)rectangle(11.9,-20.45);
\node[font=\tiny,anchor=west] at (-.4,-17.43){alt [justificatif conforme]};
\draw[<-,dashed](0,-18.25)--node[above,font=\tiny]{11a: DEFINITIVE, avertissement}(3.5,-18.25);
\draw[teal,densely dashed](-.6,-18.66)--(11.9,-18.66);
\node[font=\tiny,anchor=west] at (-.4,-18.87){[justificatif non conforme]};
\draw[<-,dashed](0,-19.76)--node[above,font=\tiny]{11b: ANNULEE, avertissement}(3.5,-19.76);
\end{tikzpicture}
\end{center}
\textbf{N1 internet -- lecture exacte :} les deux \texttt{alt} « paiement » et « justificatif » sont deux décisions \textbf{successives}. La seconde n'a lieu que si le paiement a été accepté. Le secrétariat peut saisir l'information après réception du certificat \emph{par courrier}. Le club remplace l'acteur principal et présente une \emph{licence} pour triathlon/vélo. Les réponses aux contrôles font partie du N1, même si les échanges postaux eux-mêmes ont lieu hors système.
\begin{center}\begin{tikzpicture}[x=1cm,y=.72cm]
\foreach \x/\t in {0/Concurrent,5/SystèmeCourses,10/Secrétariat}{\node[dbox] at (\x,0){\t};\draw[life](\x,-.5)--(\x,-8.6);}
\draw[assoc](0,-1)--node[above,font=\scriptsize]{1: remettreFormulaireEtPièces()}(10,-1);
\draw[assoc](10,-2)--node[above,font=\scriptsize]{2: saisirInscription()}(5,-2);
\draw[densely dashed,-Stealth](5,-3)--node[above,font=\scriptsize]{3: résultat contrôles}(10,-3);
\draw[teal] (4,-3.6) rectangle (10.75,-8.2);
\node[font=\scriptsize,anchor=west] at (4.2,-3.95){alt [valide, avant départ moins 1h et places]};
\draw[assoc](10,-4.7)--node[above,font=\scriptsize]{4a: enregistrer définitivement()}(5,-4.7);
\draw[densely dashed,-Stealth](5,-5.7)--node[above,font=\scriptsize]{5a: confirmer inscription}(10,-5.7);
\draw[teal,densely dashed](4,-6.1)--(10.75,-6.1);
\node[font=\scriptsize,anchor=west] at (4.2,-6.4){[sinon]};
\draw[densely dashed,-Stealth](5,-7.25)--node[above,font=\scriptsize]{4b: rejet et motif}(10,-7.25);
\end{tikzpicture}\end{center}
\textbf{N1 sur place :} la branche est refusée aussi pour triathlon/vélo. Pour un paiement carte, la Banque reste un acteur secondaire distinct (ajouter son échange avant confirmation).
\begin{center}\begin{tikzpicture}[x=1cm,y=.72cm]
\foreach \x/\t in {0/Concurrent,5/Secrétariat,10/SystèmeCourses}{\node[dbox] at (\x,0){\t};\draw[life](\x,-.5)--(\x,-8.1);}
\draw[assoc](0,-1)--node[above,font=\scriptsize]{1: envoyerDossierPapier()}(5,-1);
\draw[assoc](5,-2)--node[above,font=\scriptsize]{2: saisirEtContrôlerDossier()}(10,-2);
\draw[densely dashed,-Stealth](10,-3)--node[above,font=\scriptsize]{3: état du dossier}(5,-3);
\draw[teal](4.2,-3.6)rectangle(10.6,-7.8);
\node[font=\scriptsize,anchor=west] at (4.4,-3.9){alt [complet, reçu avant J--5]};
\draw[densely dashed,-Stealth](10,-4.7)--node[above,font=\scriptsize]{4a: inscription définitive}(5,-4.7);
\draw[teal,densely dashed](4.2,-5.15)--(10.6,-5.15);
\node[font=\scriptsize,anchor=west] at (4.4,-5.46){[pièce manquante ou paiement non conforme]};
\draw[densely dashed,-Stealth](10,-6.25)--node[above,font=\scriptsize]{4b: maintien état adéquat}(5,-6.25);
\draw[assoc](5,-7.05)--node[above,font=\scriptsize]{5b: demander pièce au concurrent}(0,-7.05);
\end{tikzpicture}\end{center}
\textbf{N1 courrier :} la règle J--5 s'applique à la validation finale ; une pièce manquante peut être relancée, sans annulation automatique immédiate. Pour une inscription de club, le courrier vient du Club.

\subsection{4. Classes du domaine : noms, propriétés, opérations, relations}
\begin{center}\begin{tikzpicture}[scale=.86,transform shape]
\node[dbox,align=left,minimum width=35mm] (co) at (0,3) {\textbf{Concurrent}\\nom : String\\naissance : Date\\sexe : Sexe\\estMajeurPour(e) : Bool};
\node[dbox,align=left,minimum width=37mm] (ins) at (5,3) {\textbf{Inscription}\\date : Date\\dossard : String\\statut : StatutInscription\\valider() / annuler()};
\node[dbox,align=left,minimum width=40mm] (ep) at (10,3) {\textbf{Epreuve}\\dateDépart : DateTime\\quotaMax : Integer\\tarif : Money\\estAdmissible(c) : Bool};
\node[dbox,align=left,minimum width=35mm] (club) at (0,0) {\textbf{Club}\\nom : String\\discipline : Discipline};
\node[dbox,align=left,minimum width=35mm] (pa) at (5,0) {\textbf{Paiement}\\montant : Money\\mode : ModePaiement\\estAccepté() : Bool};
\node[dbox,align=left,minimum width=37mm] (just) at (10,0) {\textbf{Justificatif}\\dateValidité : Date\\type : TypeJustif\\estConforme() : Bool};
\node[dbox,align=left,minimum width=35mm] (temps) at (0,-3) {\textbf{Passage}\\instant : DateTime\\tempsNet : Durée};
\node[dbox,align=left,minimum width=35mm] (trans) at (5,-3) {\textbf{Transpondeur}\\codeDossard : String\\estRestitué : Bool};
\node[dbox,align=left,minimum width=37mm] (cla) at (10,-3) {\textbf{Classement}\\position : Integer\\tempsFinal : Durée\\calculerPosition()};
\draw[navy](co)--node[above,font=\scriptsize,pos=.10]{1}node[above,font=\scriptsize,pos=.86]{0..*}(ins);
\draw[navy](ins)--node[above,font=\scriptsize,pos=.10]{0..*}node[above,font=\scriptsize,pos=.86]{1}(ep);
\draw[navy](co)--node[left,font=\scriptsize]{0..1 / 0..*}(club);
\draw[navy](ins)--node[left,font=\scriptsize]{1 / 0..1}(pa);
\draw[navy](ins) to[bend left=13] node[right,font=\scriptsize,pos=.5]{1 / 0..1} (just);
\draw[navy](ins)--node[above,sloped,font=\scriptsize]{0..1}(trans);
\draw[navy](trans)--node[above,font=\scriptsize]{0..*}(temps);
\draw[navy](ins)--node[above,sloped,font=\scriptsize]{0..1}(cla);
\end{tikzpicture}\end{center}
\textbf{Précision critique :} les multiplicités abrégées sur les liens obliques ci-dessus se lisent comme des \textbf{hypothèses de modélisation}, à détailler et justifier : une Inscription concerne exactement un Concurrent et une Epreuve ; un Concurrent peut avoir plusieurs Inscription ; un Paiement est rattaché au plus à une inscription individuelle dans ce modèle simplifié. Une catégorie de classement est une \textbf{valeur métier}, pas une nouvelle personne. Un Transpondeur peut être réaffecté dans le temps : les passages restent attachés à l'Inscription (ajouter une association précise en version de production).
\textbf{Classes supplémentaires utiles :} \texttt{Edition}, \texttt{Troncon}, \texttt{Antenne}, \texttt{Recompense}, \texttt{Notification}, \texttt{ClubClassement}. Ne pas stocker un rang figé avant l'arrivée : la position est calculée depuis les temps. \textbf{Classe d'association} : \texttt{Inscription} porte les attributs du lien Concurrent--Epreuve. Les contraintes d'âge et de quotas appartiennent à \texttt{Epreuve} ; les statuts à \texttt{Inscription}.

\textbf{Multiplicités à maîtriser pour ne pas perdre de points :}\par\smallskip
\begin{tabularx}{\linewidth}{@{}p{55mm}p{20mm}X@{}}
\toprule
Association (de A vers B) & A vers B & Lecture et hypothèse \\
\midrule
\texttt{Concurrent -- Inscription} & \texttt{0..* / 1} & Un concurrent peut s'inscrire plusieurs fois ; une inscription correspond à un concurrent.\\
\texttt{Epreuve -- Inscription} & \texttt{0..* / 1} & Une épreuve reçoit plusieurs inscriptions ; chacune vise une seule épreuve.\\
\texttt{Inscription -- Paiement} & \texttt{0..1 / 1} & Un dossier peut attendre le règlement ; un règlement enregistré appartient à un dossier (H2).\\
\texttt{Inscription -- Justificatif} & \texttt{0..* / 1} & Un dossier peut avoir des pièces successives ; chaque pièce est reliée à un dossier (H7).\\
\texttt{Epreuve -- Troncon} & \texttt{1..* / 1} & Parcours divisé en tronçons ; un tronçon est rattaché à son épreuve (H8).\\
\texttt{Troncon -- Antenne} & \texttt{1 / 1} & Une antenne à la fin de chaque tronçon : simplification du TD (H8).\\
\texttt{Transpondeur -- Passage} & \texttt{0..* / 1} & Plusieurs passages peuvent provenir d'une puce ; un passage porte son émetteur.\\
\texttt{Inscription -- Classement} & \texttt{0..1 / 1} & Un classement individuel est lié à une inscription terminée ; aucun rang avant résultat (H9).\\
\bottomrule
\end{tabularx}
\textbf{Sens de la table :} \texttt{A -- B : 0..* / 1} signifie qu'un A peut avoir plusieurs B et qu'un B est rattaché à un seul A. Ce sont des \textbf{choix de modélisation justifiés}, pas des multiplicités toutes écrites dans l'énoncé.\\
\textbf{Ne pas confondre :} \texttt{Inscription} rattache un \texttt{Concurrent} à une \texttt{Epreuve}. Un \texttt{Club} n'est pas un type de \texttt{Concurrent}. Le classement \textbf{par club} est calculé avec les cinq meilleurs temps ; si au plus quatre membres terminent, le club est disqualifié.

\textbf{Vue métier complémentaire -- chronométrage et récompenses :}
\begin{center}
\begin{tikzpicture}[scale=.86,transform shape]
\node[dbox] (ep) at (0,1.6) {Epreuve};
\node[dbox] (tr) at (3,1.6) {Troncon};
\node[dbox] (ant) at (6,1.6) {Antenne};
\node[dbox] (pas) at (9,1.6) {Passage};
\node[dbox] (trans) at (9,-.1) {Transpondeur};
\node[dbox] (ins) at (3,-.1) {Inscription};
\node[dbox] (clas) at (3,-1.8) {Classement};
\node[dbox] (rew) at (7.5,-1.8) {Recompense};
\draw[navy](ep)--node[above,font=\tiny,pos=.25]{1}node[above,font=\tiny,pos=.78]{1..*}(tr);
\draw[navy](tr)--node[above,font=\tiny,pos=.22]{1}node[above,font=\tiny,pos=.80]{1}(ant);
\draw[navy](ant)--node[above,font=\tiny,pos=.22]{1}node[above,font=\tiny,pos=.80]{0..*}(pas);
\draw[navy](trans)--node[right,font=\tiny]{1 / 0..*}(pas);
\draw[navy](ins)--node[left,font=\tiny]{1 / 0..1}(clas);
\draw[navy](clas)--node[above,font=\tiny]{0..* (H9)}(rew);
\end{tikzpicture}
\end{center}
\textbf{Noter H8--H9 :} un transpondeur ne remplace pas la classe \texttt{Inscription}. Pour relier chaque temps au concurrent de la bonne course, une affectation de puce à inscription est nécessaire (association ou classe d'association si l'affectation porte des dates).\par
\subsection{5. Paquetages et dépendances}
\begin{center}\begin{tikzpicture}[node distance=12mm and 20mm]
\node[cbox] (s) {\textbf{Inscriptions}\\Concurrent, Inscription, Paiement\\Justificatif};
\node[cbox,right=of s] (ev) {\textbf{Epreuves et clubs}\\Epreuve, Club, Edition};
\node[cbox,below=of s] (ch) {\textbf{Chronométrage}\\Transpondeur, Passage\\Troncon, Antenne};
\node[cbox,right=of ch] (res) {\textbf{Classements}\\Classement, Récompense};
\draw[umlarrow](s)--(ev);\draw[umlarrow](res)--(s);\draw[umlarrow](res)--(ch);\draw[umlarrow](ch)--(ev);
\end{tikzpicture}\end{center}
\textbf{Critère :} découpage par responsabilité métier ; les dépendances pointent vers le paquetage \emph{utilisé}. Pas de dépendance circulaire. Si une opération impose un cycle, déplacer un concept dans un noyau partagé ou revoir le sens du lien.

\subsection{6. États de l'inscription individuelle avec déclencheurs complets}
\begin{center}\begin{tikzpicture}[node distance=15mm and 20mm]
\node[state] (pre) {Préinscription};
\node[state,right=of pre] (fin) {Inscription financière};
\node[state,right=of fin] (def) {Inscription définitive};
\node[state,below=of fin] (ann) {Annulée};
\node[circle,fill=navy,inner sep=2pt,left=10mm of pre] (ini) {};
\node[draw,circle,minimum size=7mm,inner sep=0pt,right=14mm of def] (end) {}; \node[fill=navy,circle,minimum size=3mm,inner sep=0pt] at (end.center) {};
\draw[assoc](ini)--node[above,font=\tiny]{préinscrire [quota et date OK]}(pre);
\draw[assoc](pre)--node[above,font=\tiny]{paiementAccepté / passerFinancière()}(fin);
\draw[assoc](fin)--node[above,font=\tiny]{pièceConforme / confirmer()}(def);
\draw[assoc](def)--(end);
\draw[assoc](pre) to[bend right=28] node[below,font=\tiny]{paiementRefusé / annuler()} (ann);
\draw[assoc](fin)--node[right,font=\tiny]{certificatRefusé / annuler()}(ann);
\end{tikzpicture}\end{center}
\textbf{Énumération :} \texttt{StatutInscription = \{PREINSCRIPTION, FINANCIERE, DEFINITIVE, ANNULEE\}}. Dans le parcours \textbf{courrier}, si une pièce manque on conserve PREINSCRIPTION ou FINANCIERE et on peut reprendre le contrôle : représenter une \textbf{boucle de régularisation} \texttt{pièceReçue [conforme] / continuer()}. Dans le parcours \textbf{sur place}, trois étapes peuvent être franchies dans la même opération (pas forcément trois écrans). \textbf{Délai :} J--15 bloque l'entrée internet, J--5 la validation courrier et 1 h le guichet ; ne pas inventer un \texttt{after(24h)} inexistant dans ce TD. Le mot \texttt{after(durée)} du cours est pertinent seulement quand un délai relatif déclenche réellement un changement d'état.

\subsection{7. Deux DCP : inscription internet et recherche classement}
\begin{center}\begin{tikzpicture}[node distance=10mm and 12mm]
\node[actor] (a) {Concurrent};
\node[cbox,right=of a] (v) {«dialogue»\\PageInscription};
\node[cbox,right=of v] (c) {«contrôle»\\ControleInscription};
\node[cbox,right=of c,yshift=26mm] (e) {«entité»\\Epreuve};
\node[cbox,right=of c,yshift=13mm] (i) {«entité»\\Inscription};
\node[cbox,right=of c,yshift=0mm] (con) {«entité»\\Concurrent};
\node[cbox,right=of c,yshift=-13mm] (p) {«entité»\\Paiement};
\node[cbox,right=of c,yshift=-26mm] (j) {«entité»\\Justificatif};
\draw[navy](a)--(v);\draw[navy](v)--(c);\draw[navy](c)--(e);\draw[navy](c)--(i);\draw[navy](c)--(con);\draw[navy](c)--(p);\draw[navy](c)--(j);
\end{tikzpicture}\end{center}
\textbf{DCP inscription :} cinq entités du modèle : \texttt{Concurrent}, \texttt{Epreuve}, \texttt{Inscription}, \texttt{Paiement}, \texttt{Justificatif}. \texttt{PagePaiement} et \texttt{PageRecapitulatif} sont des dialogues \textbf{supposés} (H6), non fournis comme maquettes. La seule maquette réellement donnée est celle de recherche des classements. Le contrôle orchestre ; les entités valident la règle métier.
\begin{center}\begin{tikzpicture}[node distance=9mm and 12mm]
\node[actor] (a) {Visiteur};\node[cbox,right=of a] (v) {«dialogue»\\PageRecherche};
\node[cbox,right=of v] (c) {«contrôle»\\ControleRecherche};
\node[cbox,right=of c,yshift=16mm] (e) {«entité»\\Concurrent};
\node[cbox,right=of c] (i) {«entité»\\Inscription};
\node[cbox,right=of c,yshift=-16mm] (cl) {«entité»\\Classement};
\draw[navy](a)--(v);\draw[navy](v)--(c);\draw[navy](c)--(e);\draw[navy](c)--(i);\draw[navy](c)--(cl);
\end{tikzpicture}\end{center}
\textbf{DCP classement :} les critères exacts de la maquette sont \texttt{Dossard, Nom, Prénom, Sexe, Catégorie}. Les entités doivent reprendre exactement les noms du domaine : \texttt{Concurrent, Inscription, Classement}. \textbf{Interdit :} un lien direct Dialogue--Entité. Un \texttt{ControleRecherche} distinct est justifié car c'est un autre cas d'utilisation.

\subsection{8. Communication pour rechercher un classement}
\begin{center}\begin{tikzpicture}[node distance=15mm and 22mm]
\node[actor] (v) {Visiteur};\node[cbox,right=of v] (vue) {:PageRecherche};
\node[cbox,right=of vue] (ctrl) {:ControleRecherche};
\node[cbox,right=of ctrl] (dom) {:Classement};
\draw[navy](v)--node[above,font=\tiny]{1: saisirCritères()}(vue);
\draw[navy](vue)--node[above,font=\tiny]{2: rechercher()}node[below,font=\tiny]{5: afficher()}(ctrl);
\draw[navy](ctrl)--node[above,font=\tiny]{3: obtenirRésultats()}node[below,font=\tiny]{4: résultats}(dom);
\end{tikzpicture}\end{center}
\textbf{Complément de modèle :} \texttt{ControleRecherche} peut demander au \texttt{Concurrent} les identités et à \texttt{Inscription} les dossards. L'ordre des messages est écrit \textbf{sur les liens}, pas sur un axe du temps. Une numérotation détaillée des appels imbriqués \texttt{1, 1.1, 1.2} est aussi valide si cohérente.

\subsection{9. Séquence N3 MVC : inscription sur place}
\begin{center}\begin{tikzpicture}[x=.82cm,y=.54cm]
\foreach \x/\t in {0/Agent,3/PageGuichet,6/ControleInscription,9/Epreuve,12/Inscription}{\node[dbox,minimum width=21mm] at (\x,0){\t};\draw[life](\x,-.5)--(\x,-12.5);}
\draw[assoc](0,-1)--node[above,font=\tiny]{1: saisirDossier()}(3,-1);
\draw[assoc](3,-2)--node[above,font=\tiny]{2: inscrireSurPlace()}(6,-2);
\draw[assoc](6,-3)--node[above,font=\tiny]{3: verifierQuotaAgeDélai()}(9,-3);
\draw[densely dashed,-Stealth](9,-4)--node[above,font=\tiny]{4: résultat}(6,-4);
\draw[teal](5.25,-4.6)rectangle(12.8,-11.9);
\node[font=\tiny,anchor=west] at (5.45,-4.9){alt [admissible]};
\draw[assoc](6,-5.8)--node[above,font=\tiny]{5a: contrôlerPaiementPièce()}(12,-5.8);
\draw[densely dashed,-Stealth](12,-6.75)--node[above,font=\tiny]{6a: statut contrôle}(6,-6.75);
\draw[assoc](6,-7.7)--node[above,font=\tiny]{7a: confirmer()}(12,-7.7);
\draw[assoc](6,-8.6)--node[above,font=\tiny]{8a: afficherConfirmation()}(3,-8.6);
\draw[teal,densely dashed](5.25,-9.2)--(12.8,-9.2);
\node[font=\tiny,anchor=west] at (5.45,-9.5){[quota atteint / délai / justificatif invalide]};
\draw[assoc](6,-10.4)--node[above,font=\tiny]{5b: afficherRefus(motif)}(3,-10.4);
\end{tikzpicture}\end{center}
\textbf{Participants métier à conserver dans une version exhaustive :} \texttt{Concurrent}, \texttt{Paiement}, \texttt{Justificatif}, au besoin \texttt{Banque} pour carte et \texttt{Club} en exclusion du guichet. Le contrôle appelle les \emph{opérations métier des entités}; il ne calcule pas lui-même les seuils d'âge et de délai. L'interface ne parle jamais directement à \texttt{Epreuve} ni à \texttt{Inscription}. Le N3 reprend les interactions externes du N1.

\textbf{N3 détaillée -- messages qui doivent exister, même si le schéma précédent est compact :}
\begin{longtable}{@{}p{13mm}p{52mm}p{80mm}@{}}
\toprule
No & Émetteur vers destinataire & Message et rôle \\
\midrule\endhead
1 & Agent \(\to\) \texttt{PageGuichet} & \texttt{saisirDossier(...)}. Vue : saisie des informations et pièces.\\
2 & Vue \(\to\) \texttt{ControleInscription} & \texttt{inscrireSurPlace(...)} ; le contrôle orchestre.\\
3 & Contrôle \(\to\) \texttt{Concurrent} & \texttt{retrouverOuCreer(...)} ; réutiliser les données d'une précédente course.\\
4 & Contrôle \(\to\) \texttt{Epreuve} & \texttt{verifierAdmissibilite()}. Règles du métier : \textbf{1 h}, quota et âge ; \textbf{pas triathlon/vélo}.\\
5a & \texttt{Epreuve} \(\to\) Contrôle & \texttt{refus(motif)} ; contrôle \(\to\) vue \texttt{afficherRefus(motif)}. Fin du cas si inadmissible.\\
5b & Contrôle \(\to\) \texttt{Paiement} & \texttt{verifier(mode,montant)} ; en cas de carte, solliciter aussi la \texttt{Banque} externe.\\
6b & Contrôle \(\to\) \texttt{Justificatif} & \texttt{estConforme()} pour le certificat médical récent.\\
7b & Contrôle \(\to\) \texttt{Inscription} & \texttt{enregistrerEtValider()} seulement si \textbf{tous les contrôles sont OK}.\\
8b & Contrôle \(\to\) \texttt{PageGuichet} & \texttt{afficherConfirmation(dossard)} puis vue \(\to\) Agent.\\
\bottomrule
\end{longtable}
\textbf{Branches à écrire dans la N3 :} \texttt{alt [épreuve non admissible] / [admissible]} ; puis \texttt{alt [paiement invalide OU certificat invalide] / [tous conformes]}. En cas d'erreur, Le contrôle informe la vue ; le dossier ne devient pas \texttt{DEFINITIVE}. \textbf{H10 :} la manière de tracer une saisie incomplète n'est pas spécifiée dans le texte ; ne pas inventer une notification SMS pour le guichet.\par
\textbf{Cohérence forte :} conserver les mêmes noms de dialogues, contrôles et entités entre le DCP et la N3. Les noms \texttt{Concurrent}, \texttt{Epreuve}, \texttt{Paiement}, \texttt{Justificatif} et \texttt{Inscription} reprennent ceux du domaine.

\subsection{Dernière vérification des neuf réponses}
\begin{tabularx}{\linewidth}{@{}p{24mm}X@{}}\toprule
Questions & Critère de cohérence à vérifier avant de rendre \\\midrule
1--2 & Même acteur principal, même scénario, différences club / individuel claires.\\
3 & Trois canaux modélisés, cas d'erreur et retours, système noir.\\
4--5 & Noms métier et multiplicités ; aucune classe de DAO dans le domaine, pas de cycle.\\
6 & Énumération alignée et branches d'annulation/régularisation présentes.\\
7--8 & Dialogue--Contrôle--Entités, acteurs à l'extérieur, mêmes noms, messages numérotés.\\
9 & N3 MVC détaillée, mêmes événements que N1, contrôle court, règles dans les entités.\\\bottomrule\end{tabularx}
\clearpage
\section{Planches complémentaires : solutions lisibles pour l'examen}
Les diagrammes suivants \textbf{complètent} ceux des neuf réponses précédentes. Ils sont volontairement plus grands et séparés par scénario. Chaque diagramme indique les messages, les alternatives et les sorties d'erreur, afin de servir de modèle de dessin sur papier.

\subsection{Q1 -- Vue complète des acteurs et fonctionnalités}
\begin{center}\begin{tikzpicture}[x=1cm,y=1cm,scale=.87,transform shape]
\node[actor] (public) at (0,3.8) {Public};
\node[actor] (conc) at (0,2.3) {Concurrent};
\node[actor] (club) at (0,.8) {Club};
\node[actor] (sec) at (0,-1.1) {Secrétariat};
\node[actor] (sc) at (0,-3.2) {Service courses};
\node[actor] (bank) at (9.3,2.5) {Banque};
\node[actor] (chrono) at (9.3,-2.4) {Chronométrage};
\node[actor] (sms) at (9.3,-3.8) {Service SMS};
\draw[teal,rounded corners=3pt,thick] (2.1,-5) rectangle (7.9,5.0);
\node[font=\scriptsize,fill=white] at (5,5){Application de gestion des courses};
\node[usecase] (info) at (4.8,4.1) {Consulter\\conditions};
\node[usecase] (form) at (4.8,2.95) {Télécharger\\formulaire};
\node[usecase] (internet) at (4.8,1.8) {S'inscrire en ligne};
\node[usecase] (papier) at (4.8,.65) {Enregistrer\\courrier};
\node[usecase] (guichet) at (4.8,-.5) {Inscrire sur place};
\node[usecase] (controle) at (4.8,-1.65) {Contrôler\\paiement et pièces};
\node[usecase] (result) at (4.8,-2.8) {Gérer résultats\\et récompenses};
\node[usecase] (recherche) at (4.8,-3.95) {Rechercher\\classement};
\draw[navy] (public)--(info);\draw[navy] (public)--(form);\draw[navy] (public) to[bend right=25] (recherche);
\draw[navy] (conc)--(internet);\draw[navy] (club)--(internet);\draw[navy] (club)--(papier);
\draw[navy] (sec)--(papier);\draw[navy] (sec)--(guichet);\draw[navy] (sec)--(controle);
\draw[navy] (sc)--(result);\draw[navy] (bank)--(internet);\draw[navy] (chrono)--(result);\draw[navy] (sms)--(result);
\end{tikzpicture}\end{center}
\textbf{Lecture.} Les acteurs à droite représentent ici des systèmes externes (H2). Ces acteurs échangent avec l'application. Le transpondeur seul n'est pas automatiquement un acteur : c'est le système de chronométrage qui peut jouer ce rôle. Un club ne doit pas être associé à « inscrire sur place ». Pour préciser la variante par club, une seconde vue de cas (voir question 1) évite de surcharger le schéma.

\subsection{Q2 -- Description textuelle : les quatre sorties à ne pas confondre}
\begin{tabularx}{\linewidth}{@{}p{27mm}X@{}}\toprule
Sortie & Formulation française à utiliser dans la copie \\ \midrule
Succès & \textit{« L'inscription devient définitive quand le paiement et le justificatif sont acceptés. »}\\
Préinscription refusée & \textit{« Si le quota est atteint ou la date J--15 dépassée, aucune préinscription n'est enregistrée. »}\\
Paiement refusé & \textit{« Le dossier est annulé et le concurrent est informé. »}\\
Justificatif refusé & \textit{« Le secrétariat refuse la pièce ; le système annule le dossier et avertit le concurrent. »}\\
\bottomrule\end{tabularx}
\textbf{Question piège :} « Précondition = il reste des places » ? Les places sont une condition métier contrôlée dans le déroulement ; si elles manquent, le cas doit expliquer le refus. Il est plus précis de réserver les préconditions aux éléments nécessaires avant le déclenchement (accès au site et épreuve visible) puis d'écrire les vérifications dans le scénario.
\newpage
\subsection{Q3 -- N1 Internet : décision de préinscription, paiement, pièce}
\begin{center}\begin{tikzpicture}[x=1cm,y=.61cm]
\foreach \x/\t in {0/Concurrent,3.3/Système,6.6/Banque,9.8/Secrétariat}{\node[dbox,minimum width=24mm] at (\x,0){\t};\draw[life](\x,-.6)--(\x,-17.2);}
\draw[assoc](0,-1.3)--node[above,font=\scriptsize]{1: choisirEpreuve()}(3.3,-1.3);
\draw[densely dashed,-Stealth](3.3,-2.3)--node[above,font=\scriptsize]{2: conditions, tarif, quota}(0,-2.3);
\draw[assoc](0,-3.4)--node[above,font=\scriptsize]{3: saisirInformations()}(3.3,-3.4);
\draw[teal](-.5,-3.9)rectangle(10.55,-6.9);
\node[font=\scriptsize,anchor=west] at(-.3,-4.2){alt [âge, quota et J--15 conformes]};
\draw[densely dashed,-Stealth](3.3,-5.2)--node[above,font=\tiny]{4a: préinscription créée}(0,-5.2);
\draw[teal,densely dashed](-.5,-5.6)--(10.55,-5.6);
\node[font=\scriptsize,anchor=west] at(-.3,-5.9){[non conforme]};
\draw[densely dashed,-Stealth](3.3,-6.35)--node[above,font=\tiny]{4b: refus motivé}(0,-6.35);
\draw[assoc](0,-7.7)--node[above,font=\scriptsize]{5: payerParCarte()}(3.3,-7.7);
\draw[assoc](3.3,-8.8)--node[above,font=\scriptsize]{6: demanderAccord()}(6.6,-8.8);
\draw[densely dashed,-Stealth](6.6,-9.8)--node[above,font=\scriptsize]{7: accord / refus}(3.3,-9.8);
\draw[teal](-.5,-10.2)rectangle(10.55,-13.1);
\node[font=\scriptsize,anchor=west] at(-.3,-10.5){alt [paiement accepté]};
\draw[densely dashed,-Stealth](3.3,-11.5)--node[above,font=\tiny]{8a: état FINANCIERE}(0,-11.5);
\draw[teal,densely dashed](-.5,-11.9)--(10.55,-11.9);
\node[font=\scriptsize,anchor=west] at(-.3,-12.2){[paiement refusé]};
\draw[densely dashed,-Stealth](3.3,-12.65)--node[above,font=\tiny]{8b: ANNULEE + information}(0,-12.65);
\draw[assoc](9.8,-14.1)--node[above,font=\scriptsize]{9: enregistrerContrôlePièce()}(3.3,-14.1);
\draw[teal](-.5,-14.5)rectangle(10.55,-17.0);
\node[font=\scriptsize,anchor=west] at(-.3,-14.8){alt [certificat conforme]};
\draw[densely dashed,-Stealth](3.3,-15.65)--node[above,font=\tiny]{10a: DEFINITIVE}(0,-15.65);
\draw[teal,densely dashed](-.5,-16.05)--(10.55,-16.05);
\node[font=\scriptsize,anchor=west] at(-.3,-16.35){[certificat refusé]};
\draw[densely dashed,-Stealth](3.3,-16.75)--node[above,font=\tiny]{10b: ANNULEE + information}(0,-16.75);
\end{tikzpicture}\end{center}
\begin{warningbox}
L'exemple regroupe \textbf{trois \texttt{alt} successifs}. Dans une séquence UML stricte, arrêter le parcours après les branches \texttt{refus} (ou utiliser un \texttt{break}) : une annulation n'est \textbf{jamais} suivie du paiement ou du contrôle d'un certificat. Le courrier du certificat est hors du système web ; seul son contrôle enregistré génère un message système. Pour le paiement par chèque, remplacer l'échange bancaire par celui du secrétariat.
\end{warningbox}
\medskip
\subsection{Q3 -- N1 Sur place : décision et refus sans triathlon/vélo}
\begin{center}\begin{tikzpicture}[x=1cm,y=.7cm]
\foreach \x/\t in {0/Concurrent,4.7/Secrétariat,9.4/Système}{\node[dbox,minimum width=27mm] at(\x,0){\t};\draw[life](\x,-.5)--(\x,-12.2);}
\draw[assoc](0,-1.2)--node[above,font=\scriptsize]{1: remettreDossierEtPaiement()}(4.7,-1.2);
\draw[assoc](4.7,-2.3)--node[above,font=\scriptsize]{2: saisirDonnéesEtJustificatif()}(9.4,-2.3);
\draw[densely dashed,-Stealth](9.4,-3.4)--node[above,font=\scriptsize]{3: admissibilité + statut}(4.7,-3.4);
\draw[teal](3.55,-4.2) rectangle(10.1,-11.7);
\node[font=\scriptsize,anchor=west] at(3.8,-4.6){alt [places, âge, mode sur place, 1h OK]};
\draw[assoc](4.7,-5.6)--node[above,font=\tiny]{4a: confirmer Paiement, Certificat()}(9.4,-5.6);
\draw[densely dashed,-Stealth](9.4,-6.6)--node[above,font=\tiny]{5a: inscription définitive}(4.7,-6.6);
\draw[assoc](4.7,-7.6)--node[above,font=\tiny]{6a: remettreConfirmation()}(0,-7.6);
\draw[teal,densely dashed](3.55,-8.4)--(10.1,-8.4);
\node[font=\scriptsize,anchor=west] at(3.8,-8.8){[quota/âge/heure/mode/paiement/pièce KO]};
\draw[densely dashed,-Stealth](9.4,-9.8)--node[above,font=\tiny]{4b: refus motivé}(4.7,-9.8);
\draw[assoc](4.7,-10.9)--node[above,font=\tiny]{5b: informerRefus()}(0,-10.9);
\end{tikzpicture}\end{center}
\textbf{Particularité :} les 3 étapes sont réalisées en une seule visite. Le statut final est \texttt{DEFINITIVE} uniquement si tous les contrôles passent. Pas de prise d'inscription sur place pour les triathlons et courses de vélo. La banque intervient pour la carte : détailler sa demande et sa réponse si le paiement carte est choisi. Cette planche reprend la vue principale en noir, sans tous les détails du prestataire externe.
\medskip
\subsection{Q3 -- N1 Courrier : un dossier peut attendre une régularisation}
\begin{center}\begin{tikzpicture}[x=1cm,y=.66cm]
\foreach \x/\t in {0/Concurrent,4.7/Secrétariat,9.4/Système}{\node[dbox,minimum width=27mm] at(\x,0){\t};\draw[life](\x,-.5)--(\x,-13.4);}
\draw[assoc](0,-1.3)--node[above,font=\scriptsize]{1: envoyerDossierPapier()}(4.7,-1.3);
\draw[assoc](4.7,-2.5)--node[above,font=\scriptsize]{2: saisirDossierEtPièces()}(9.4,-2.5);
\draw[densely dashed,-Stealth](9.4,-3.6)--node[above,font=\scriptsize]{3: retour contrôles}(4.7,-3.6);
\draw[teal](3.6,-4.2) rectangle(10.1,-12.9);
\node[font=\scriptsize,anchor=west] at(3.85,-4.55){alt [dossier complet et date <= J--5]};
\draw[assoc](4.7,-5.6)--node[above,font=\tiny]{4a: confirmerDossier()}(9.4,-5.6);
\draw[densely dashed,-Stealth](9.4,-6.7)--node[above,font=\tiny]{5a: DEFINITIVE}(4.7,-6.7);
\draw[teal,densely dashed](3.6,-7.5)--(10.1,-7.5);
\node[font=\scriptsize,anchor=west] at(3.85,-7.85){[pièce ou paiement non conforme]};
\draw[densely dashed,-Stealth](9.4,-8.85)--node[above,font=\tiny]{4b: statut conservé}(4.7,-8.85);
\draw[assoc](4.7,-9.9)--node[above,font=\tiny]{5b: demanderDocument()}(0,-9.9);
\draw[assoc](0,-11.0)--node[above,font=\tiny]{6b: fournirComplément()}(4.7,-11.0);
\draw[assoc](4.7,-12.05)--node[above,font=\tiny]{7b: recontrôlerDossier()}(9.4,-12.05);
\end{tikzpicture}\end{center}
\textbf{Important :} la branche de complément \emph{réessaie} seulement tant qu'elle respecte J--5. Le texte ne dit pas qu'une simple pièce manquante entraîne l'annulation immédiate. On conserve la phase déjà atteinte et on contacte le concurrent.
\medskip
\subsection{Q3 -- N1 Club : ce qui change vraiment}
\begin{center}\begin{tikzpicture}[x=1cm,y=.72cm]
\foreach \x/\t in {0/Club,4.7/Système,9.4/Secrétariat}{\node[dbox,minimum width=27mm] at(\x,0){\t};\draw[life](\x,-.5)--(\x,-10.0);}
\draw[assoc](0,-1.2)--node[above,font=\scriptsize]{1: inscrireMembre(membre, épreuve)}(4.7,-1.2);
\draw[densely dashed,-Stealth](4.7,-2.3)--node[above,font=\scriptsize]{2: résultat contrôles club et licence}(0,-2.3);
\draw[teal](-.55,-3.1)rectangle(10.1,-9.6);
\node[font=\scriptsize,anchor=west] at(-.3,-3.4){alt [club spécialisé, places, J--15, membre licencié]};
\draw[assoc](0,-4.5)--node[above,font=\tiny]{3a: demanderInscription()}(4.7,-4.5);
\draw[assoc](9.4,-5.6)--node[above,font=\tiny]{4a: confirmerPaiementEtLicence()}(4.7,-5.6);
\draw[densely dashed,-Stealth](4.7,-6.7)--node[above,font=\tiny]{5a: DEFINITIVE}(0,-6.7);
\draw[teal,densely dashed](-.55,-7.4)--(10.1,-7.4);
\node[font=\scriptsize,anchor=west] at(-.3,-7.75){[club / licence / délai non conformes]};
\draw[densely dashed,-Stealth](4.7,-8.75)--node[above,font=\tiny]{3b: refus motivé}(0,-8.75);
\end{tikzpicture}\end{center}
\textbf{Hypothèse :} le contrôle documentaire se fait par le secrétariat, comme pour les autres épreuves. Le scénario du club \textbf{reprend la validation du paiement}, mais le certificat médical est \textbf{remplacé par la licence} : ne dessinez pas les deux documents comme obligatoires en même temps.
\newpage
\subsection{Q4 -- Classes du domaine : tableaux de multiplicité sans ambiguïté}
\small\begin{longtable}{@{}p{42mm}p{13mm}p{14mm}p{80mm}@{}}\toprule
Association (A -- B) & vers B & vers A & Lecture / justification \\ \midrule\endhead
Concurrent -- Inscription & 0..* & 1 & Un concurrent peut faire plusieurs inscriptions ; chaque inscription vise une personne.\\
Epreuve -- Inscription & 0..* & 1 & Une course reçoit plusieurs dossiers ; un dossier concerne une course.\\
Inscription -- Paiement & 0..1 & 1 & Modèle simplifié : un seul paiement associé au dossier, ou aucun si en attente. H2.\\
Inscription -- Justificatif & 0..* & 1 & Plusieurs pièces successives possibles (nouveau certificat) ; chacune appartient à un dossier. H7.\\
Club -- Concurrent & 0..* & 0..1 & Un membre peut être inscrit pour un club dans ce modèle simplifié ; affiliation à un instant donné. H5.\\
Epreuve -- Troncon & 1..* & 1 & Un parcours est découpé en tronçons ; chaque tronçon appartient à une épreuve. H8.\\
Troncon -- Antenne & 1 & 1 & Antenne de fin de tronçon ; simplification pédagogique (la technologie réelle peut réutiliser des équipements). H8.\\
Transpondeur -- Passage & 0..* & 1 & Une puce déclenche plusieurs temps de passage ; un passage vient d'une puce.\\
Inscription -- Classement & 0..1 & 1 & Le résultat final du coureur n'existe qu'après la course ; le dossier peut ne pas avoir de classement. H9.\\
Club -- ClassementClub & 0..* & 1 & Plusieurs participations de club au fil des compétitions ; chaque classement de club vise un club. H9.\\
\bottomrule\end{longtable}\normalsize
\textbf{Vérification utile :} les multiplicités se placent \emph{près de l'extrémité} qu'elles quantifient. Les tables ci-dessus utilisent « vers B » pour la quantité de B par instance de A. L'association du club doit pouvoir représenter l'historique si les adhésions changent ; une classe d'association \texttt{Adhesion} serait alors préférable (H5).

\subsection{Q4 -- Contrat de données pour les entités les plus utilisées}
\small\begin{longtable}{@{}p{29mm}p{63mm}p{61mm}@{}}\toprule
Entité & Attributs minimum (exemples) & Opérations métier possibles \\ \midrule\endhead
Concurrent & id, nom, prénom, dateNaissance, sexe, téléphone & verifierAgePour(e), actualiserCoordonnees()\\
Epreuve & id, type, dateHeureDepart, quota, tarif & placesDisponibles(), accepteInscription(c,date)\\
Inscription & id, date, dossard, statut, canal & preinscrire(), enregistrerPaiement(), valider(), annuler()\\
Paiement & montant, mode, date, accepte & controlerValidite(), enregistrerAccord()\\
Justificatif & type, dateEmission, conforme & estConformePour(e), renouvelerControle()\\
Club & id, nom, discipline & estSpecialisePour(e), inscrireMembre()\\
Transpondeur & code, dossardCourant, restitue & attribuer(), marquerRestitue()\\
Passage & dateHeure, tempsIntermediaire & calculerTempsNet()\\
Classement & sexe, categorie, position, temps & calculerPosition(), publier()\\
ClassementClub & totalTop5, rang, disqualifie & calculerCinqMeilleurs(), disqualifierSiMoinsDeCinq()\\
\bottomrule\end{longtable}\normalsize
\textit{Précision :} « opérations possibles » signifie conception de responsabilités cohérente avec le texte, pas signatures imposées par l'enseignant. Le domaine ne contient pas \texttt{SQLConnection}, \texttt{RepositoryDAO}, \texttt{VueHTML} ou \texttt{SessionHTTP}.
\medskip
\subsection{Q6 -- Etats : événements, gardes, actions, y compris régularisation}
\begin{center}\begin{tikzpicture}[node distance=12mm and 18mm]
\node[state] (pre) {Préinscription};
\node[state,right=of pre] (fin) {Financière};
\node[state,right=of fin] (def) {Définitive};
\node[state,below=24mm of fin] (ann) {Annulée};
\node[circle,fill=navy,inner sep=2pt,left=10mm of pre] (start) {};
\node[draw,circle,minimum size=6mm,inner sep=0pt,right=13mm of def] (end) {};
\node[fill=navy,circle,minimum size=2.7mm,inner sep=0pt] at(end.center){};
\draw[assoc](start)--node[above,font=\tiny]{demande [quota, âge, date OK] / créer()}(pre);
\draw[assoc](pre)--node[above,font=\tiny]{paiementAccepté [accord obtenu] / noterAccord()}(fin);
\draw[assoc](fin)--node[above,font=\tiny]{pièceConforme [moins de 1 an OU licence] / confirmer()}(def);
\draw[assoc](def)--(end);
\draw[assoc](pre) to[bend right=23] node[below,sloped,font=\tiny]{refusPaiement [canal internet] / annuler()}(ann);
\draw[assoc](fin)--node[right,font=\tiny]{pièceRefusée [canal internet] / annuler()}(ann);
\draw[assoc](pre) to[out=130,in=230,looseness=5] node[left,font=\tiny]{pièceReçue [courrier, avant J--5] / contrôler()} (pre);
\draw[assoc](fin) to[out=35,in=145,looseness=5] node[above,font=\tiny]{pièceReçue [courrier, avant J--5] / vérifier()}(fin);
\end{tikzpicture}\end{center}
\textbf{Lecture importante :} le modèle est \emph{une vue abstraite}. Sur place ou avec un dossier postal complet, le logiciel peut traverser les trois états en une seule opération : \texttt{Préinscription} puis \texttt{Financière} puis \texttt{Définitive}. Le délai J--5 est une garde pour le courrier, pas un \texttt{after(5 jours)} si aucune annulation automatique n'est donnée. L'état \texttt{ANNULEE} est terminal dans ce modèle, sauf nouvelle demande distincte.
\newpage
\subsection{Q7 -- DCP : deux cas, deux contrôles (schémas détaillés)}
\begin{center}\begin{tikzpicture}[node distance=9mm and 12mm]
\node[actor] (a) {Concurrent};\node[cbox,right=of a] (d1) {«dialogue»\\PageSaisie\\saisirDonnees()};
\node[cbox,below=of d1] (d2) {«dialogue»\\PagePaiement\\choisirMode()};
\node[cbox,right=18mm of d1] (ctrl) {«contrôle»\\ControleInscription\\demanderInscription()};
\node[cbox,right=20mm of ctrl,yshift=20mm] (e) {«entité»\\Epreuve};
\node[cbox,right=20mm of ctrl,yshift=10mm] (c) {«entité»\\Concurrent};
\node[cbox,right=20mm of ctrl] (i) {«entité»\\Inscription};
\node[cbox,right=20mm of ctrl,yshift=-10mm] (p) {«entité»\\Paiement};
\node[cbox,right=20mm of ctrl,yshift=-20mm] (j) {«entité»\\Justificatif};
\draw[navy](a)--(d1);\draw[navy](a)--(d2);\draw[navy](d1)--(ctrl);\draw[navy](d2)--(ctrl);
\foreach \target in {e,c,i,p,j}{\draw[navy](ctrl)--(\target);}
\end{tikzpicture}\end{center}
\textbf{H6 :} les écrans de saisie et de paiement sont déduits du scénario ; \emph{ils ne sont pas fournis} comme maquettes dans le TD. Une vraie maquette n'existe que pour « Rechercher un classement ».
\begin{center}\begin{tikzpicture}[node distance=10mm and 16mm]
\node[actor] (v) {Public};
\node[cbox,right=of v] (vue) {«dialogue»\\PageRechercheClassement\\saisirCritères()};
\node[cbox,right=of vue] (ctrl) {«contrôle»\\ControleRecherche\\rechercher()};
\node[cbox,right=of ctrl,yshift=13mm] (a) {«entité»\\Concurrent};
\node[cbox,right=of ctrl] (b) {«entité»\\Inscription};
\node[cbox,right=of ctrl,yshift=-13mm] (c) {«entité»\\Classement};
\draw[navy](v)--(vue);\draw[navy](vue)--(ctrl);\draw[navy](ctrl)--(a);\draw[navy](ctrl)--(b);\draw[navy](ctrl)--(c);
\end{tikzpicture}\end{center}
\textbf{Règle de traçabilité :} mêmes noms d'entités dans le diagramme de classes et le DCP. La vue ne doit pas appeler l'entité directement. Un contrôle est lié à un cas d'utilisation, pas à chaque bouton.
\medskip
\subsection{Q8 -- Communication : la chronologie est écrite sur les liens}
\begin{center}\begin{tikzpicture}[node distance=14mm and 20mm]
\node[actor] (act) {Visiteur};\node[cbox,right=of act] (vue) {:PageRechercheClassement};
\node[cbox,right=of vue] (ctrl) {:ControleRecherche};
\node[cbox,right=of ctrl,yshift=15mm] (c) {:Concurrent};
\node[cbox,right=of ctrl] (i) {:Inscription};
\node[cbox,right=of ctrl,yshift=-15mm] (cl) {:Classement};
\draw[navy](act)--node[above,font=\tiny]{1: saisirCritères()}(vue);
\draw[navy](vue)--node[above,font=\tiny]{2: rechercher(critères)}node[below,font=\tiny]{6: afficher(résultats)}(ctrl);
\draw[navy](ctrl)--node[above,sloped,font=\tiny]{3: lireIdentités()}(c);
\draw[navy](ctrl)--node[above,font=\tiny]{4: lireDossards()}(i);
\draw[navy](ctrl)--node[below,sloped,font=\tiny]{5: trouverClassement()}(cl);
\end{tikzpicture}\end{center}
\textbf{Si aucun classement n'est trouvé :} message conditionnel \texttt{6 [résultat vide] : afficherAucunRésultat()}. Les éléments \texttt{Dossard, Nom, Prénom, Sexe, Catégorie} viennent de la capture de l'énoncé (page 5). Le diagramme de communication numérote les messages \textbf{sur les liens} ; il n'a pas de temps vertical comme la séquence.
\newpage
\subsection{Q9 -- Séquence N3 MVC, inscription sur place, partie 1}
\begin{center}\begin{tikzpicture}[x=1cm,y=.65cm]
\foreach \x/\t in {0/Secrétariat,3.2/PageGuichet,6.5/ControleInscription,9.8/Epreuve}{\node[dbox,minimum width=25mm] at(\x,0){\t};\draw[life](\x,-.6)--(\x,-12.3);}
\draw[assoc](0,-1.2)--node[above,font=\scriptsize]{1: saisirDossier()}(3.2,-1.2);
\draw[assoc](3.2,-2.3)--node[above,font=\scriptsize]{2: inscrireSurPlace()}(6.5,-2.3);
\draw[assoc](6.5,-3.4)--node[above,font=\scriptsize]{3: estAdmissible()}(9.8,-3.4);
\draw[densely dashed,-Stealth](9.8,-4.45)--node[above,font=\scriptsize]{4: places, âge, délai, canal}(6.5,-4.45);
\draw[teal](2.6,-5.1) rectangle(10.4,-11.8);
\node[font=\scriptsize,anchor=west] at(2.8,-5.45){alt [non admissible]};
\draw[assoc](6.5,-6.55)--node[above,font=\tiny]{5a: afficherRefus(motif)}(3.2,-6.55);
\draw[teal,densely dashed](2.6,-7.5)--(10.4,-7.5);
\node[font=\scriptsize,anchor=west] at(2.8,-7.82){[admissible]};
\draw[assoc](6.5,-8.95)--node[above,font=\tiny]{5b: ouvrirContrôlePièces()}(3.2,-8.95);
\draw[densely dashed,-Stealth](3.2,-10.05)--node[above,font=\tiny]{6b: demanderPièces()}(0,-10.05);
\end{tikzpicture}\end{center}
\textbf{Attention :} ce schéma traite seulement l'admissibilité. Ensuite, l'entité \texttt{Inscription} peut créer le dossier. La seconde planche montre les contrôles suivants.
\medskip
\subsection{Q9 -- Séquence N3 MVC, inscription sur place, partie 2}
\begin{center}\begin{tikzpicture}[x=.9cm,y=.62cm]
\foreach \x/\t in {0/Agent,3/Vue,6/Contrôle,9/Paiement,12/Justificatif,15/Inscription}{\node[dbox,minimum width=17mm,font=\tiny] at(\x,0){\t};\draw[life](\x,-.5)--(\x,-14.1);}
\draw[assoc](0,-1.15)--node[above,font=\tiny]{1: transmettreDonnées()}(3,-1.15);
\draw[assoc](3,-2.25)--node[above,font=\tiny]{2: enregistrerPièces()}(6,-2.25);
\draw[assoc](6,-3.35)--node[above,font=\tiny]{3: verifier(mode,montant)}(9,-3.35);
\draw[densely dashed,-Stealth](9,-4.45)--node[above,font=\tiny]{4: accord / refus}(6,-4.45);
\draw[assoc](6,-5.55)--node[above,font=\tiny]{5: estConforme()}(12,-5.55);
\draw[densely dashed,-Stealth](12,-6.65)--node[above,font=\tiny]{6: conforme / refus}(6,-6.65);
\draw[teal](5.4,-7.2) rectangle(15.65,-13.7);
\node[font=\scriptsize,anchor=west] at(5.6,-7.55){alt [tous les contrôles sont OK]};
\draw[assoc](6,-8.65)--node[above,font=\tiny]{7a: validerInscription()}(15,-8.65);
\draw[densely dashed,-Stealth](15,-9.75)--node[above,font=\tiny]{8a: confirmation + dossard}(6,-9.75);
\draw[assoc](6,-10.85)--node[above,font=\tiny]{9a: afficherConfirmation()}(3,-10.85);
\draw[teal,densely dashed](5.4,-11.35)--(15.65,-11.35);
\node[font=\scriptsize,anchor=west] at(5.6,-11.65){[paiement ou certificat refusé]};
\draw[assoc](6,-12.8)--node[above,font=\tiny]{7b: afficherRefus()}(3,-12.8);
\end{tikzpicture}\end{center}
\textbf{Lecture MVC :} l'agent parle à la vue ; la vue parle au contrôle ; le contrôle appelle des entités métier. L'entité \texttt{Paiement} ne demande pas elle-même un affichage ; \texttt{Justificatif} décide de sa conformité. \textbf{Brancher la banque} lorsque le mode carte l'exige. \textbf{Avant cette partie}, le contrôle a aussi récupéré/actualisé \texttt{Concurrent} et vérifié \texttt{Epreuve} : voir tableau détaillé de la question 9.
\newpage
\section{Contrôle qualité avant impression et avant l'examen}
\begin{tabularx}{\linewidth}{@{}p{19mm}X@{}}\toprule
Travail & Ce qui doit apparaître dans votre copie \\ \midrule
Q1 & Acteurs du domaine et systèmes externes, cas cohérents, périmètre, variante club interdite au guichet.\\
Q2 & Identification, but, déclencheur, préconditions, chemin nominal, alternatives, exceptions, postconditions, délais et règlement.\\
Q3 & Internet, guichet, courrier, club en variante ; N1 = système en boîte noire ; aucune classe technique.\\
Q4 & Classes métier, opérations vraies, données, associations et multiplicités justifiées ; chronométrage et récompenses.\\
Q5 & Paquetages orientés métier, dépendances orientées vers celui qui est utilisé.\\
Q6 & État initial et final, 4 statuts, déclencheurs \texttt{événement [garde] / action}, reprise courrier.\\
Q7 & Deux DCP ; acteurs -- dialogues -- contrôles -- entités, aucun lien direct dialogue--entité.\\
Q8 & Objets identifiés, liens et messages numérotés ; branche de classement non trouvé.\\
Q9 & Même acteurs et décisions qu'en N1 ; vue, contrôle et entités cohérents avec DCP, paiement et certificat.\\
\bottomrule\end{tabularx}
\begin{warningbox}
Ne confondez pas \textbf{correction détaillée} et \textbf{vérité du sujet}. L'énoncé ne définit pas exactement toutes les classes, cardinalités, écrans internes ni le contenu des notifications. Les parties déduites doivent être écrites comme \textbf{hypothèses} et non présentées comme faits. C'est une compétence demandée dans ASI C5.
\end{warningbox}
\section*{Sources}
\addcontentsline{toc}{section}{Sources}
\begin{itemize}
\item \textit{Organisation de courses de fond}, énoncé universitaire, pages 3--6 : règles métier et neuf questions.
\item ASI C1--C3 : cas d'utilisation, description textuelle, classes, états, DCP, N1 et N3.
\item ASI C5 : fiches des six diagrammes et contrôle de cohérence.
\item Copie personnelle envoyée à M. Ali Lahmer, le 21 octobre 2026 : point de comparaison, \textbf{pas} preuve d'une correction officielle.
\end{itemize}
\end{document}